\chapter[Band 2: Theoreme der Logik]{Band 2 auf einen Blick}
\noindent\textbf{Theoreme der Logik}\par\medskip
\noindent\emph{Lesart.} \(P,Q,R\) sind beliebige Formeln;
\(P(x)\) und \(Q(x)\) dürfen freie Variablen besitzen. Die
Grundjunktoren und Quantoren stammen aus Band 1; dieser Band beweist
ihre Rechengesetze und definiert weitere logische Operationen.
\begin{BandUebersicht}
\textbf{Konjunktion und Disjunktion}
Die iterierte Schreibweise wird rechts geklammert:
\UebFormel{\begin{aligned}
P\land Q\land R&\coloneqq P\land(Q\land R),\\
P\lor Q\lor R&\coloneqq P\lor(Q\lor R).
\end{aligned}}
\UebQuellen{\FormulaRefAuto{P \land Q \land R \coloneqq P \land (Q \land R)}[def];
\FormulaRefAuto{P \lor Q \lor R \coloneqq P \lor (Q \lor R)}[def]}
&
\textit{Idempotenz, Kommutativität und Distribution}
\UebFormel{\begin{aligned}
&P\land P\eqvdash P,\quad P\lor P\eqvdash P,\\
&P\land Q\eqvdash Q\land P,\\
&P\land(Q\lor R)\\
&\hspace{8mm}\eqvdash(P\land Q)\lor(P\land R).
\end{aligned}}
\UebQuellen{\FormulaRefAuto{P \land P \eqvdash P}[thm];
\FormulaRefAuto{P \lor P \eqvdash P}[thm];
\FormulaRefAuto{P\land Q\eqvdash Q\land P}[thm];
\FormulaRefAuto{P \land (Q\lor R) \eqvdash (P \land Q) \lor (P\land R)}[thm]}\\
\addlinespace[.8ex]
\textbf{Negation und Wahrheitswerte}
\UebFormel{\top\coloneqq\neg\bot.}
\(\bot\) und \(\top\) erlauben die Formulierung der Wahrheitstafeln
als beweisbare Äquivalenzen.
\UebQuellen{\FormulaRefAuto{\top \coloneqq \neg\bot}[def]}
&
\textit{Klassische Negationsgesetze}
\UebFormel{\begin{aligned}
&P\eqvdash\neg\neg P,\qquad P\lor\neg P,\\
&\neg(P\lor Q)\eqvdash\neg P\land\neg Q,\\
&\neg(P\land Q)\eqvdash\neg P\lor\neg Q,\\
&\bot\vdash P.
\end{aligned}}
\UebQuellen{\FormulaRefAuto{P \eqvdash \neg\neg P}[thm];
\FormulaRefAuto{P \lor \neg P}[thm];
\FormulaRefAuto{\neg(P \lor Q) \eqvdash \neg P \land \neg Q}[thm];
\FormulaRefAuto{\neg(P \land Q) \eqvdash \neg P \lor \neg Q}[thm];
\FormulaRefAuto{\bot\vdash P}[thm]}\\
\addlinespace[.8ex]
\textbf{Implikation und Bikonditional}
Die Grundzeichen \(\rightarrow\) und \(\leftrightarrow\) werden
mittels ihrer Einführungs- und Eliminationsregeln verwendet.
&
\textit{Charakterisierungen und Kontraposition}
\UebFormel{\begin{aligned}
&P\rightarrow Q\eqvdash\neg P\lor Q,\\
&P\rightarrow Q\eqvdash\neg Q\rightarrow\neg P,\\
&P\leftrightarrow Q\\
&\hspace{8mm}\eqvdash(P\rightarrow Q)\land(Q\rightarrow P),\\
&(P\land Q)\rightarrow R\\
&\hspace{8mm}\eqvdash P\rightarrow(Q\rightarrow R).
\end{aligned}}
\UebQuellen{\FormulaRefAuto{P \rightarrow Q \eqvdash \neg P \lor Q}[thm];
\FormulaRefAuto{P \rightarrow Q \eqvdash \neg Q \rightarrow \neg P}[thm];
\FormulaRefAuto{P \leftrightarrow Q \eqvdash (P \rightarrow Q)\land (Q\rightarrow P)}[thm];
\FormulaRefAuto{(P \land Q) \rightarrow R \eqvdash P \rightarrow (Q \rightarrow R)}[thm]}\\
\addlinespace[.8ex]
\textbf{Exklusives Oder}
\UebFormel{P\lxor Q\coloneqq\neg(P\leftrightarrow Q).}
\UebQuellen{\FormulaRefAuto{P \lxor Q \coloneqq \neg(P \leftrightarrow Q)}[def]}
&
\textit{Genau einer der beiden Fälle}
\UebFormel{\begin{aligned}
&P\lxor Q\eqvdash(\neg P\land Q)\lor(P\land\neg Q),\\
&P\lxor Q\eqvdash Q\lxor P,\qquad\neg(P\lxor P).
\end{aligned}}
\UebQuellen{\FormulaRefAuto{P \lxor Q \eqvdash (\neg P\land Q)\lor(P\land \neg Q)}[thm];
\FormulaRefAuto{P\lxor Q \eqvdash Q\lxor P}[thm];
\FormulaRefAuto{\neg(P\lxor P)}[thm]}\\
\addlinespace[.8ex]
\textbf{NAND und NOR}
\UebFormel{\begin{aligned}
P\lnand Q&\coloneqq\neg(P\land Q),\\
P\lnor Q&\coloneqq\neg(P\lor Q).
\end{aligned}}
\UebQuellen{\FormulaRefAuto{P \lnand Q \coloneqq \neg(P \land Q)}[def];
\FormulaRefAuto{P \lnor Q \coloneqq \neg(P \lor Q)}[def]}
&
\textit{Darstellung anderer Junktoren}
\UebFormel{\begin{aligned}
&P\lnand P\eqvdash\neg P,\\
&P\land Q\eqvdash(P\lnand Q)\lnand(P\lnand Q),\\
&P\lnor P\eqvdash\neg P,\\
&P\lor Q\eqvdash(P\lnor Q)\lnor(P\lnor Q).
\end{aligned}}
\UebQuellen{\FormulaRefAuto{P\lnand P \eqvdash \neg P}[thm];
\FormulaRefAuto{P\land Q \eqvdash (P\lnand Q)\lnand(P\lnand Q)}[thm];
\FormulaRefAuto{P\lnor P \eqvdash \neg P}[thm];
\FormulaRefAuto{P\lor Q \eqvdash (P\lnor Q)\lnor(P\lnor Q)}[thm]}\\
\addlinespace[.8ex]
\textbf{Halbaddierer}
\UebFormel{\begin{aligned}
\HAddSum{P}{Q}&\coloneqq P\lxor Q,\\
\HAddCarry{P}{Q}&\coloneqq P\land Q.
\end{aligned}}
\UebQuellen{\FormulaRefAuto{\HAddSum{P}{Q}\coloneqq P\lxor Q}[def];
\FormulaRefAuto{\HAddCarry{P}{Q}\coloneqq P\land Q}[def]}
&
\textit{Summe und Übertrag schließen einander aus}
\UebFormel{\begin{aligned}
&\HAddSum{P}{Q}\vdash\neg\HAddCarry{P}{Q},\\
&P\dsep Q\vdash
\neg\HAddSum{P}{Q}\land\HAddCarry{P}{Q}.
\end{aligned}}
\UebQuellen{\FormulaRefAuto{\HAddSum{P}{Q}\vdash \neg\HAddCarry{P}{Q}}[thm];
\FormulaRefAuto{P\dsep Q\vdash \neg\HAddSum{P}{Q}\land \HAddCarry{P}{Q}}[thm]}\\
\addlinespace[.8ex]
\textbf{Volladdierer}
\UebFormel{\begin{aligned}
&\FAddSum{P}{Q}{R}\coloneqq\\
&\qquad\HAddSum{\HAddSum{P}{Q}}{R},\\
&\FAddCarry{P}{Q}{R}\coloneqq\HAddCarry{P}{Q}\\
&\qquad\lor\HAddCarry{\HAddSum{P}{Q}}{R}.
\end{aligned}}
\UebQuellen{\FormulaRefAuto{\FAddSum{P}{Q}{R}\coloneqq \HAddSum{\HAddSum{P}{Q}}{R}}[def];
\FormulaRefAuto{\FAddCarry{P}{Q}{R}\coloneqq \HAddCarry{P}{Q}\lor \HAddCarry{\HAddSum{P}{Q}}{R}}[def]}
&
\textit{Explizite Formeln}
\UebFormel{\begin{aligned}
&\FAddSum{P}{Q}{R}\eqvdash(P\lxor Q)\lxor R,\\
&\FAddCarry{P}{Q}{R}\\
&\hspace{6mm}\eqvdash(P\land Q)\lor((P\lxor Q)\land R).
\end{aligned}}
\UebQuellen{\FormulaRefAuto{\FAddSum{P}{Q}{R}\eqvdash (P\lxor Q)\lxor R}[thm];
\FormulaRefAuto{\FAddCarry{P}{Q}{R}\eqvdash (P\land Q)\lor((P\lxor Q)\land R)}[thm]}\\
\addlinespace[.8ex]
\textbf{Identität und Ungleichheit}
Die Gleichheit besitzt die Regeln aus Band 1; \(a\neq b\)
ist die Negation von \(a=b\).
&
\textit{Kongruenz von Funktionstermen}
\UebFormel{\begin{aligned}
&a=b\vdash f(c,a)=f(c,b),\\
&a=b\vdash f(a,c)=f(b,c).
\end{aligned}}
\UebQuellen{\FormulaRefAuto{BinaryFunctionRightArgumentCongruence}[thm];
\FormulaRefAuto{BinaryFunctionLeftArgumentCongruence}[thm]}\\
\addlinespace[.8ex]
\textbf{Quantoren und Eindeutigkeit}
\UebFormel{\begin{aligned}
&\existsleone x\,P(x)\coloneqq\\
&\quad\forall y\forall z\,((P(y)\land P(z))\rightarrow y=z).
\end{aligned}}
\UebQuellen{\FormulaRefAuto{\existsleone x\,P(x) \coloneqq \forall y\forall z((P(y)\land P(z))\rightarrow y=z)}[def]}
&
\textit{Quantorennegation und eindeutige Existenz}
\UebFormel{\begin{aligned}
&\neg\forall x\,P(x)\eqvdash\exists x\,\neg P(x),\\
&\forall x\,\neg P(x)\eqvdash\neg\exists x\,P(x),\\
&\exists!x\,P(x)\eqvdash
\exists x\,P(x)\land\existsleone x\,P(x).
\end{aligned}}
\UebQuellen{\FormulaRefAuto{\neg\forall x(P(x)) \eqvdash \exists x (\neg P(x))}[thm];
\FormulaRefAuto{\forall x(\neg P(x)) \eqvdash \neg\exists x (P(x))}[thm];
\FormulaRefAuto{\exists! x\,P(x) \eqvdash \exists x\,P(x)\land \existsleone x\,P(x)}[thm]}\\
\addlinespace[.8ex]
\textbf{Konditionaler Aussagenoperator und Fallterm}
\UebFormel{\begin{aligned}
\IfWff{P}{Q}{R}&\coloneqq(P\land Q)\lor(\neg P\land R),\\
&\IfThenElse{P}{a}{b}.
\end{aligned}}
Der Fallterm bezeichnet \(a\) im Fall \(P\), andernfalls \(b\).
\UebQuellen{\FormulaRefAuto{if-}[def];\FormulaRefAuto{Fallterm}[def]}
&
\textit{Implikationsform und Auswertung}
\UebFormel{\begin{aligned}
&\IfWff{P}{Q}{R}\\
&\hspace{6mm}\eqvdash(P\rightarrow Q)\land(\neg P\rightarrow R),\\
&P\vdash\IfThenElse{P}{a}{b}=a,\\
&\neg P\vdash\IfThenElse{P}{a}{b}=b.
\end{aligned}}
\UebQuellen{\FormulaRefAuto{if-Implikationsform}[thm];
\FormulaRefAuto{P \vdash \IfThenElse{P}{a}{b}=a}[thm];
\FormulaRefAuto{\neg P \vdash \IfThenElse{P}{a}{b}=b}[thm]}\\
\end{BandUebersicht}
